\section{Knowledge-Guided Imputation (KGI) Across Architectures}

The integration of Domain Knowledge into time-series imputation models represents a core contribution of this work. We propose Knowledge-Guided Imputation (KGI), a framework designed to infuse statistical deep learning models with medical domain knowledge extracted from the Unified Medical Language System (UMLS). Depending on the underlying model's architecture, KGI is injected either through structural relational priors (Graph Structure Learning) or dynamic contextual text embeddings (Semantic Fusion). Here, we detail the distinct injection mechanisms across the four primary baseline architectures evaluated in our study.

\subsection{SAITS: Multi-Modal Semantic Fusion and Graph Priors}
The Self-Attention-based Imputation for Time Series (SAITS) architecture incorporates KGI at two fundamental levels: semantic cross-attention and graph-structured priors.

\subsubsection*{Semantic Injection via Multimodal Fusion}
During the early forward pass, the base latent tensor $\mathbf{H}_{base} \in \mathbb{R}^{B \times T \times d_{model}}$ is intercepted before the primary DMSA blocks. A \texttt{TextualKnowledgeInjector} dynamically extracts pre-computed MedBERT embeddings representing clinical relational triplets based on the patient's currently observed variables. A Multi-Head Cross-Attention layer is then employed where the numerical latent tensor $\mathbf{H}_{base}$ serves as the Query ($\mathbf{Q}$), while the medical semantic tensor $\mathbf{C}_{text}$ acts as both Key ($\mathbf{K}$) and Value ($\mathbf{V}$). The semantically enriched output is added back via a residual connection, allowing the numerical hidden state to be modulated by medical language rules before any intra-series attention is computed.

\subsubsection*{Relational Graph Prior Initialization}
Alternatively, SAITS accommodates structural KGI using SapBERT embeddings. A static relational prior matrix $\mathbf{P}$ is computed using cosine similarity among the variables' textual embeddings. The attention maps within the \texttt{GraphDMSABlock} are initialized and regularized against this physiological prior $\mathbf{P}$. A graph structure learning loss $||\mathbf{A} - \mathbf{P}||_F$ ensures the learned temporal dependencies do not deviate excessively from established medical knowledge.

\subsection{BRITS: Contextual Recurrent Injection}
Bidirectional Recurrent Imputation for Time Series (BRITS) relies on a dual-LSTM bidirectional architecture. To inject KGI, we fundamentally alter the recurrent state updates within the \texttt{BackboneRITS} modules for both the forward and backward directions.

At each time step $t$, the standard BRITS model concatenates the complementarily imputed feature vector $\mathbf{c}_c$ with the missing indicator mask $\mathbf{m}$ to feed the LSTM cell. In our KGI-enhanced formulation, the \texttt{TextualKnowledgeInjector} provides a text-derived contextual vector $\mathbf{C}_{text}^{(t)}$ that encapsulates the semantic meaning of the currently observed variables. The LSTM input dimension is explicitly expanded to accept the concatenation $[\mathbf{c}_c, \mathbf{m}, \mathbf{C}_{text}^{(t)}]$. By injecting the knowledge graph semantics directly into the recurrent block, the hidden states simultaneously evolve over time driven by both numerical regressions and medical linguistic rules.

\subsection{MRNN: Fully Connected Network Refinement}
The Multi-directional Recurrent Neural Network (MRNN) separates the imputation process into two stages: an initial sequence estimation via bidirectional RNNs, followed by a Fully Connected Network (FCN) regression that refines the prediction using cross-feature correlations.

We target the refinement phase (\texttt{MrnnFcnRegression}) to inject KGI. The RNN produces an intermediate hidden representation $\mathbf{h}_t$, which in the vanilla MRNN is directly projected to the final imputed value $\mathbf{\hat{x}}_t$ via a sigmoid layer. In our KGI adaptation, the temporal semantic context vector $\mathbf{C}_{text}$ is retrieved based on the patient's observable dataset mask. We concatenate the intermediate numerical representation with these text embeddings: $[\mathbf{h}_t, \mathbf{C}_{text}]$. A new linear layer processes this joint text-numerical representation to yield the final imputation. This design ensures that while the RNN captures raw temporal dynamic features, the FCN leverages broad medical semantics to correct unrealistic clinical physiological outliers prior to finalizing the prediction.

\subsection{TimesFM: Graph Feature Interaction via Patching}
TimesFM operates primarily on sequence patches rather than individual time steps, mapping well to foundation model paradigms. KGI is implemented in this architecture as an explicit topological message-passing layer named \texttt{GraphFeatureInteraction}.

Before processing patches through the transformer backbone, we introduce a trainable adjacency matrix $\mathbf{A} \in \mathbb{R}^{D \times D}$, representing the interactions between the $D$ clinical features. $\mathbf{A}$ is explicitly initialized using the SapBERT relational prior $\mathbf{P}$. The layer performs message passing across features at the patch-token level, computing $\mathbf{H}_{interact} = \text{Softmax}(\mathbf{A}) \times \mathbf{H}_{patch}$. To preserve numerical stability and prevent signal collapse, this interaction is integrated back into the main flow via a learnable multi-layer gating mechanism and a residual connection: $\mathbf{H}_{out} = \mathbf{H}_{patch} + \sigma(\mathbf{Gate}) \odot \mathbf{H}_{interact}$. This approach structurally forces the TimesFM token representations to exchange information strictly along predefined physiological graph hierarchies right from the input layer.
